%-----------------------------------------------------------------------
%                          hires_correction_letter.tex
% A&A LaTeX class (aa.cls)
%-----------------------------------------------------------------------
%\documentclass[referee]{aa}    % for a referee version
%\documentclass[letter]{aa}     % for the letters
\documentclass{aa}
\usepackage{graphicx}
\usepackage{txfonts}
\usepackage{placeins}
\usepackage{natbib}
\usepackage[colorlinks,allcolors=blue,bookmarks=false]{hyperref}
\usepackage{xcolor}
\usepackage{amsmath}
\hyphenation{get-sources get-filaments get-images non-uniform back-ground}

\newcommand{\draftnote}[1]{\textcolor{magenta}{\textbf{[#1]}}}
\newcommand{\chg}[1]{\textcolor{magenta}{#1}}
\newcommand{\mH}{m_{\rm H}}
\newcommand{\muH}{\mu_{\rm H_2}}
\newcommand{\um}{\,$\mu$m}
\newcommand{\Teff}{T_{\rm eff}}
\newcommand{\Sc}{\mathcal{D}_{\rm c}}
\newcommand{\Sim}{\mathcal{D}}
\newcommand{\Iim}{\mathcal{I}}

\begin{document}

\title{Correcting the dust temperature bias of high-resolution surface
       density images of molecular clouds}

\author{Alexander Men'shchikov\inst{1}}
\institute{D\'epartement d'Astrophysique (DAp), IRFU, CEA Saclay,
Universit\'e Paris-Saclay, 91191 Gif-sur-Yvette, France}

\date{Received / Accepted}
\offprints{Alexander Men'shchikov}

\abstract{Dust surface densities of molecular clouds are derived by fitting a
single modified blackbody to the far-infrared intensities of each pixel. A line
of sight crossing a dense structure is not isothermal, and such a fit is
weighted towards the warmer material, so the surface densities come out too low.
The bias grows with the surface density of the environment, with the
concentration of a structure, and with the degree to which it is unresolved, so
it distorts the relative surface densities of structures in different
environments and at different distances. We present a method that removes it,
without extracting sources and without calibrating on models. The image is
decomposed into levels by a granulometry, a nested set of greyscale
morphological openings; the levels are non-negative, sum exactly to the image,
and each holds material within a range of widths. The same decomposition of the
waveband images yields a dust temperature for each level, of material of that
width alone. Each level is then converted with its own
temperature rather than with the single temperature of the pixel, which
multiplies it by the ratio of the Planck functions at the two temperatures. The
correction is exactly unity where a line of sight is isothermal. On a sub-field of
the Aquila cloud it leaves the diffuse emission unchanged and raises the peaks
of dense structures by a median factor of 1.7; on radiative transfer models of
embedded cores, the median recovered fraction of the true peak rises from 0.63
to 0.84.}

\keywords{ISM: clouds -- ISM: structure -- Stars: formation --
Infrared: ISM -- Submillimeter: ISM -- Techniques: image processing}

\maketitle

% =====================================================================
\section{Introduction}
\label{sec:intro}

Far-infrared and submillimeter surveys of molecular clouds with the
\emph{Herschel} Space Observatory have produced the images from which most of
what is known about the structure of nearby star-forming regions has been
derived \citep[e.g.][]{Andre+2010, Konyves+2015, Andre+2014}. The physical
quantity of interest, the surface density of dust and gas, is not observed
directly. It is obtained by fitting an optically thin modified blackbody to the
intensities observed in each pixel \citep{Hildebrand1983}, which yields one
dust temperature and one surface density per line of sight.

The angular resolution of such an image is that of the longest waveband used,
$36\farcs3$ at 500\um, unless a differential scheme transfers to it the higher
resolution of the shorter wavebands \citep{Palmeirim+2013}. The implementation
used here, \emph{hires}, is described in
\citet[][Appendix~A]{Menshchikov2021method} and produces surface density images
at the resolution of any of the observed wavebands. Such images are the basis
of the extraction of dense cores and filaments
\citep{Menshchikov2021, Konyves+2015}, and every quantity derived from them, in
particular the core mass function, inherits their accuracy.

The single-temperature assumption is the weak point of the procedure. It is
accurate for diffuse material, whose temperature varies little along the line
of sight, and it fails for a cold dense structure embedded in a warmer cloud.
The fitted temperature is weighted by the emission, hence by the warmer
material, and the surface density derived from it is too low. The magnitude of
the effect follows from the logarithmic sensitivity of the Planck function,
%
\begin{equation}
   \frac{\partial \ln \Sigma}{\partial \ln T} =
   -\,\frac{x_\lambda}{1-{\rm e}^{-x_\lambda}}\,,
   \qquad x_\lambda \equiv \frac{h c}{\lambda k T}\,,
   \label{eq:sensitivity}
\end{equation}
%
which at $T = 10$\,K equals $-8.99$ at 160\um, $-5.78$ at 250\um, $-4.18$ at
350\um, and $-3.05$ at 500\um. Overestimating the dust temperature of a cold
core by only 1\,K therefore lowers the surface density derived from the
500\um\ intensity by 30\,per cent, and that derived from the 160\um\ intensity
by a factor of 2.3. A modest error in the temperature is a large error in the
surface density.

The bias is not a multiplicative constant that could be absorbed into the
assumed dust opacity. It grows with the surface density of the surrounding
cloud, because a denser environment produces a larger temperature contrast
along the line of sight; it grows with the concentration of the structure,
because a more centrally condensed core is more strongly shielded; and it grows
as the structure becomes unresolved, because the beam then mixes the cold
interior with warmer material across the sky as well as along the line of
sight. Two structures of the same true surface density, observed in different
environments or at different distances, are therefore recovered with different
values, and their surface densities are not comparable.

In this Letter we describe a correction that removes the bias directly from a
derived surface density image. Section~\ref{sec:method} presents the method,
Sect.~\ref{sec:aquila} illustrates it on a sub-field of the Aquila cloud,
Sect.~\ref{sec:tests} presents the tests that establish its behaviour, and
Sect.~\ref{sec:limits} states its limitations. Figure~\ref{fig:flowchart}
summarises the procedure end to end.

Throughout this paper images are denoted by capital calligraphic characters
(for instance $\Iim$, $\Sim$, $\mathcal{L}$), to distinguish them from the
scalar quantities that appear alongside them.

\begin{figure*}
  \centering
  \includegraphics[width=0.92\linewidth]{fig_method_flowchart.pdf}
  \caption{Overview of the method. The surface density image to be corrected
  (\emph{top}) and the observed waveband images brought to a common beam
  (\emph{bottom}) are decomposed into the same levels by a nested set of
  morphological openings (Sect.~\ref{sec:levels}). The image and the longest
  waveband give the effective temperature $\Teff$ with which the image was
  built, and the levels of the waveband images give a dust temperature $T_k$ for
  the material of each width (Sect.~\ref{sec:leveltemp}). Each level of the
  image is then multiplied by the ratio of the Planck functions at those two
  temperatures, $\Teff$ in the numerator and $T_k$ in the denominator, and the
  levels are summed (Sect.~\ref{sec:correction}). No source is detected at any
  stage and no calibration enters.}
  \label{fig:flowchart}
\end{figure*}

% =====================================================================
\section{Method}
\label{sec:method}

\subsection{Decomposition into levels}
\label{sec:levels}

The correction must act on material of different widths separately, because
material of different widths has different temperatures. The decomposition that
separates it is a granulometry, the classical tool of mathematical morphology
for measuring size, introduced by \citet{Matheron1975} for binary images and
extended to greyscale images by \citet{Serra1982} and \citet{Sternberg1986}.

The two operations from which it is built are simple to state. With a flat
disk $B$ of diameter $d$ as the structuring element, the erosion and the
dilation of an image $\Iim$ replace each pixel by, respectively, the smallest and
the largest value of the image within the disk centred on that pixel,
%
\begin{equation}
   [\varepsilon_d \Iim](x) = \min_{y \in B} \Iim(x+y) \,,
   \qquad
   [\delta_d \Iim](x) = \max_{y \in B} \Iim(x+y) \,,
   \label{eq:erodil}
\end{equation}
%
and the opening is the one followed by the other,
$O_d(\Iim) = \delta_d\,[\varepsilon_d \Iim]$.

The composition of the two has a geometrical meaning that is more useful than
the formula. Think of the image as a surface whose height is the intensity, and
of the disk as a rigid flat plate of diameter $d$ held horizontally. The erosion
gives, at each position, the highest level at which the plate can be held
centred there without any part of it emerging above the surface. The dilation
then spreads each of those admissible positions back over the whole area of the
plate. The result, at every pixel, is the highest level that any plate lying
entirely beneath the surface reaches at that pixel:
%
\begin{equation}
   [O_d \Iim](x) = \max_{\{B\,\ni\,x\}}\ \min_{y \in B} \Iim(y) \,.
   \label{eq:opening}
\end{equation}
%
This answers a question that the picture of a ball rolled beneath a surface
leaves open. The plate touches the surface only at isolated points, and those
contact points do not by themselves define a surface; it is the envelope of all
admissible positions of the plate, and not the set of contacts, that is the
opening, and Eq.~(\ref{eq:opening}) is defined at every pixel whether or not a
plate touches the surface there. Fluctuations narrower than $d$ are therefore
removed regardless of where they sit: a spike cannot support a plate, and a
dip on the flank of a broad structure is simply bridged by a plate resting on
either side of it.

Three properties follow directly from Eq.~(\ref{eq:opening}) and are used
below. The opening never exceeds the image, $O_d(\Iim) \le I$, since a plate lies
beneath the surface by construction; it is decreasing in the diameter,
$O_{d'}(\Iim) \le O_d(\Iim)$ for $d' > d$, since a larger plate is admissible in
fewer positions; and applying it twice changes nothing. A structure narrower
than the plate cannot hold it up and is therefore absent from $O_d(\Iim)$, whereas
a structure wider than the plate supports it and is retained. The opening is
thus a filter in width, and $d$ is a width threshold and nothing else.

The single-opening version of this operation, the top-hat transform
$I - O_d(\Iim)$, is the background subtraction used routinely in biomedical
imaging, where the structuring element is often taken to be a ball rather than
a flat plate and the operation is known as rolling-ball background subtraction
\citep{Sternberg1983}. We use a flat disk, for which
Eq.~(\ref{eq:opening}) holds exactly as written.

A granulometry applies a sequence of such openings with structuring elements of
increasing size, in the manner of sieving a granular sample through sieves of
decreasing hole size, and records what is removed at each step. With an
increasing set of diameters $d_1 < d_2 < \ldots < d_N$ we therefore define the
levels
%
\begin{equation}
   \mathcal{L}_k \equiv O_{d_{k-1}}(\Iim) - O_{d_k}(\Iim)\,,
   \qquad O_{d_0} \equiv I \,,
   \label{eq:levels}
\end{equation}
%
so that the image is recovered exactly as
%
\begin{equation}
   \Iim = O_{d_N}(\Iim) + \sum_{k=1}^{N} \mathcal{L}_k \,.
   \label{eq:sum}
\end{equation}
%
The collection $\{\mathcal{L}_k\}$ is the greyscale pattern spectrum of
\citet{Maragos1989}, whose integrals over the image give the distribution of
material over width. Here we do not use those integrals: we keep the levels
themselves as images and treat each of them as a separate physical component.
In astronomy the same construction has been used to measure the size
distribution of ionized regions in simulations of reionization
\citep{Kakiichi+2017}, on binary fields rather than on greyscale images.

Because $O_{d_k}(\Iim) \le O_{d_{k-1}}(\Iim)$ everywhere, every level is
non-negative. Material appears in the level that matches its own width, and the
widest background $O_{d_N}(\Iim)$, the diffuse cirrus, is set aside and left
uncorrected. Figure~\ref{fig:levelscut} shows the construction along a cut
across the Aquila sub-field.

Nothing is classified as a source or as a filament at any point. A structure is
assigned to a level by its width, which is a well defined quantity, whereas
deciding whether an elongated structure is a filament or a wide source is not
\citep{Menshchikov2021}. A wide source and a filament of the same width
receive the same treatment, which is what the physics requires: both are cold
material seen through warmer material, and the bias that affects them has the
same origin.

A granulometry is deliberately preferred here to a decomposition into spatial
scales. A scale decomposition, formed from differences between successive
convolutions of the image, produces components that change sign, so that a
compact source is accompanied by a negative ring and is spread over several
components. A temperature derived from such a component would be a ratio of
quantities that pass through zero, and would be correspondingly unstable; we
found its pixel-to-pixel scatter to be several kelvin, comparable to its whole
range. The levels of Eq.~(\ref{eq:levels}) are non-negative differences between
an image and a surface lying beneath it, and their temperatures are stable to
$0.2$\,K between adjacent pixels (Sect.~\ref{sec:stability}).

\begin{figure}
  \centering
  \includegraphics[width=\linewidth]{fig_levels_cut.pdf}
  \caption{Decomposition into levels along a cut across the Aquila sub-field.
  The black curve is the surface density image and the coloured curves are the
  openings $O_d$ with plates of increasing diameter, each of which passes
  beneath every structure narrower than itself. The three grey bars are
  admissible positions of the plate of diameter $91\arcsec$: each lies entirely
  beneath the surface, and the opening $O_{91\arcsec}$ is the envelope of all
  such positions (Eq.~\ref{eq:opening}), not the set of points at which a plate
  touches the surface. A level is the material between two successive curves,
  and the lowest curve is the cirrus, which is left uncorrected. The levels are
  non-negative and sum exactly to the image (Eq.~\ref{eq:sum}).}
  \label{fig:levelscut}
\end{figure}

\subsection{Dust temperature of a level}
\label{sec:leveltemp}

The same decomposition, with the same diameters, is applied to each observed
waveband image after all of them have been convolved to a common beam, that of
the longest waveband. The level $\mathcal{L}_k^\lambda$ then contains the emission of the
material of the corresponding width, while the warmer material in front of it
and behind it has been assigned to the wider levels and to the background.

Fitting an optically thin modified blackbody with a fixed emissivity index to
the set $\{\mathcal{L}_k^\lambda\}$ of one level yields a dust temperature $T_k$ that
describes material of that width alone, and not the whole line of sight. This
is the essential measurement of the method, and it is possible only because the
levels are non-negative: the spectral energy distribution of a level is a set
of positive intensities of comparable magnitude, not a set of small differences
between large numbers.

\subsection{Correction}
\label{sec:correction}

Let $\Sim$ be the surface density image to be corrected and $\lambda_{\rm r}$
the longest waveband used to build it. The image was constructed as
%
\begin{equation}
   \Sim = \frac{\Iim_{\lambda_{\rm r}}}
                 {B_{\lambda_{\rm r}}(T)\,\kappa_{\lambda_{\rm r}}\,
                  \eta\,\muH\,\mH} \,,
   \label{eq:sigma}
\end{equation}
%
where $\kappa_\lambda$ is the dust opacity, $\eta$ the dust-to-gas mass ratio,
and $\muH$ the mean molecular weight per hydrogen molecule. Its effective dust
temperature $\Teff$ is therefore recovered exactly by inverting
Eq.~(\ref{eq:sigma}), without refitting anything,
%
\begin{equation}
   B_{\lambda_{\rm r}}(\Teff) =
   \frac{I_{\lambda_{\rm r}}}
        {\Sim\,\kappa_{\lambda_{\rm r}}\,\eta\,\muH\,\mH} \,.
   \label{eq:teff}
\end{equation}
%
Every level of $\Sim$ was converted into surface density with $\Teff$,
whereas it should have been converted with its own temperature $T_k$. The two
conversions differ by the ratio of the Planck functions, so the corrected image
is
%
\begin{equation}
   \Sc = O_{d_N}(\Sim) + \sum_{k=1}^{N} C_k\, \mathcal{L}_k(\Sim) \,,
   \qquad
   C_k = \frac{B_{\lambda_{\rm r}}(\Teff)}{B_{\lambda_{\rm r}}(T_k)} \,.
   \label{eq:correction}
\end{equation}
%
No empirical calibration and no model enter Eq.~(\ref{eq:correction}). Three
properties follow immediately and are worth stating separately.

First, the correction vanishes identically where it should. If the line of
sight is isothermal, every level shares the temperature of the pixel, $T_k =
\Teff$, hence $C_k = 1$ and the image is returned unchanged. The correction
acts only where material of different widths has different temperatures, which
is precisely the condition under which the single-temperature fit is in error.
This is not an approximate statement but an algebraic identity, and it is
verified numerically in Sect.~\ref{sec:nulltest}.

Second, the corrected image can neither become negative nor develop a
depression. The levels are non-negative and each is multiplied by a factor
$C_k \ge 1$, so nothing is ever subtracted.

Third, the shape of a structure is preserved. Multiplying a level by a factor
changes its amplitude and not its profile, so a round source stays round and no
structure is deformed. This is a practical requirement rather than an aesthetic
one: an image that is to be searched for sources must not contain artefacts
that a source-extraction algorithm would mistake for structures.

The dust temperature corresponding to the corrected image follows by inverting
Eq.~(\ref{eq:teff}) with $\Sc$ in place of $\Sim$, which is the same as
dividing the Planck function by the correction factor. Defined that way it is
exactly consistent with $\Sc$ and identical to $\Teff$ wherever the correction
is unity, so that no feature can appear in it where nothing has been corrected.

\subsection{Practical considerations}
\label{sec:practical}

Five points matter in an application to observed images.

\subsubsection{The ladder of diameters}
\label{sec:ladder}

The diameters need not be chosen individually. A geometric ladder from twice
the angular resolution to a largest ball, beyond which material is treated as
cirrus and left unchanged, removes the arbitrariness. On the Aquila sub-field of
Sect.~\ref{sec:aquila} the median correction changes by 2\,per cent when the
ratio between successive diameters is varied from 1.3 to 2.0, so the result
does not depend on the choice.

\subsubsection{Coverage}
\label{sec:coverage}

The coverage masks of all wavebands are multiplied together and contracted by a
few pixels. Pixels outside the common coverage are filled with the value of the
nearest covered pixel before any opening or convolution, so that neither
operation is affected by the boundary, and the correction is tapered to unity
near the edge, because a structure cut by the edge of the coverage is not
measured correctly. Pixels in which the input image is not positive are
excluded, since inverting Eq.~(\ref{eq:teff}) there would return meaningless
temperatures.

\subsubsection{Angular resolution}
\label{sec:resolution}

The correction is a temperature effect and has no structure below the angular
resolution of the image, which has none either. The correction field
$\Sc/\Sim$ is therefore smoothed to the resolution before it is applied.
Without this step the correction, which is built from the levels, carries
structure down to the pixel scale and amplifies the noise differentially.

\subsubsection{Noise}
\label{sec:noise}

Instrumental noise is independent between wavebands, whereas real emission
appears at the same position in all of them. Regressing each waveband of a
level on the longest waveband and taking the robust scatter of the residual
therefore measures the noise of that level, in the units of the reference
waveband, without any noise image being supplied. Where a level does not stand
above its own incoherent part, the correction is faded out. There is a
systematic reason for doing so and not merely accepting the scatter: $C_k$ is
convex in $T_k$ and is bounded below by unity, so random scatter in $T_k$ would
produce a mean correction above unity even where the true answer is one.

\subsubsection{Internally heated sources}
\label{sec:protostars}

A protostar has a steep temperature gradient along the line of sight, from tens
of kelvin in its interior to about 10\,K in its envelope, and a single
temperature per level does not describe it. Such objects are identified
unambiguously by their compact emission at 70\um, and the correction is
switched off smoothly over the heated region. The 70\um\ image is used here
only as a flag and never as a photometric measurement: on the Aquila sub-field
its compact component amounts to two parts in a thousand of the local
background, with a pixel-to-pixel scatter far larger, so it carries no usable
temperature information.

% =====================================================================
\section{Application to the Aquila cloud}
\label{sec:aquila}

We illustrate the method on a sub-field of $31\arcmin \times 31\arcmin$ of the
Aquila cloud observed with \emph{Herschel} at 70, 160, 250, 350, and 500\um\
\citep{Andre+2010, Konyves+2015}, using the surface density image at
$13\farcs5$ resolution produced with \emph{hires}. The ladder of ball diameters
runs geometrically from $27\arcsec$ to $205\arcsec$ with a ratio of 1.5.
\draftnote{Adopted distance and its reference to be inserted.}

Figure~\ref{fig:leveltemp} shows the central measurement. The dust temperature
of a level rises systematically with its width, from $8.7$\,K for material
narrower than $27\arcsec$ to $10.4$\,K for material narrower than $205\arcsec$,
while the single-temperature fit of the same lines of sight gives about 15\,K.
The correction factors of Eq.~(\ref{eq:correction}) fall correspondingly, from
about 2.4 for the narrowest level to about 1.7 for the widest. That the narrow
material is colder than the wide material is not assumed anywhere; it is
measured, and it is the entire content of the correction.

\begin{figure}
  \centering
  \includegraphics[width=\linewidth]{fig_level_temperatures.pdf}
  \caption{Dust temperature of each level (black, left axis) as a function of
  the ball diameter that defines it, on the Aquila sub-field, with the median
  and the 16th to 84th percentile range over the pixels in which the level
  carries material. The dashed line is the median effective temperature $\Teff$
  of the same lines of sight, obtained from the single-temperature fit. The
  purple symbols and the shaded band (right axis) are the resulting correction
  factors $C_k$ of Eq.~(\ref{eq:correction}).}
  \label{fig:leveltemp}
\end{figure}

Figure~\ref{fig:maps} shows the original image, the corrected image, and the
correction factor. The diffuse emission is untouched and the dense structures
are raised. Over the whole sub-field the ratio $\Sc/\Sim$ has a median of
1.02, a 90th percentile of 1.35, and a maximum of 1.83, while in the faintest
30\,per cent of the image its median is 1.00. The correction factor traces the
dense material and reaches its largest values on the coldest structures, so
that the third panel of Fig.~\ref{fig:maps} is, in effect, a map of where the
single-temperature assumption fails.

\begin{figure*}
  \centering
  \includegraphics[width=\linewidth]{fig_aquila_maps.pdf}
  \caption{The Aquila sub-field, $31\arcmin \times 31\arcmin$.
  (\emph{left}) the surface density image at $13\farcs5$ produced with
  \emph{hires}; (\emph{middle}) the same image after the correction of
  Eq.~(\ref{eq:correction}), on the same scale; (\emph{right}) the correction
  factor $\Sc/\Sim$, which is unity in the diffuse emission and reaches 1.8
  on the coldest dense structures.}
  \label{fig:maps}
\end{figure*}

Measured at the positions of the 144 brightest local maxima of the original
image, the background-subtracted peak surface density rises by a median factor
of 1.71, with a 16th to 84th percentile range of 1.06 to 2.15.
Figure~\ref{fig:profile} shows the effect on the radial profile of the
brightest structure of the field: the amplitude increases while the shape is
preserved, as Sect.~\ref{sec:correction} requires.

\begin{figure}
  \centering
  \includegraphics[width=\linewidth]{fig_core_profile.pdf}
  \caption{Radial profile of the brightest structure of the Aquila sub-field,
  before and after the correction. The amplitude is raised and the shape is
  unchanged, since each level is multiplied by a factor and not reshaped.}
  \label{fig:profile}
\end{figure}

The corrected image is smooth. The absolute difference between adjacent pixels,
relative to the local value, has a median of 1.6\,per cent against 1.4\,per
cent for the original image; no pixel is negative and none is lowered.

% =====================================================================
\section{Tests}
\label{sec:tests}

\subsection{An isothermal sky}
\label{sec:nulltest}

The most informative test of Eq.~(\ref{eq:correction}) is that it must do
nothing where there is nothing to do. A simulated image of a cloud containing
cirrus, a filament, dense cores, and noise was constructed with a single dust
temperature assigned to each line of sight, and the waveband images were
computed from it. In such an image the single-temperature fit is exact by
construction and no correction is due.

The method returns factors below 1.01 for the cirrus and for the filament, that
is, it leaves them alone as it must. The dense cores of the same simulation
were computed from radiative transfer models and therefore do possess a
temperature structure along the line of sight; they are corrected by up to a
factor of 2.1. The method thus distinguishes, within one image and with no
external information, between cold material that is biased and cold material
that is not.

\subsection{Stability of the level temperatures}
\label{sec:stability}

On the Aquila sub-field the dust temperature of a level has a root mean square
difference of $0.2$\,K between adjacent pixels, against a spread of $3.7$\,K
over the field. The ratio of the two, about one to twenty, is the property that
makes the method usable: the quantity that drives the correction is measured
far more precisely than it varies.

The median temperature of material narrower than $60\arcsec$ is $9.9$\,K on the
Aquila sub-field and $9.7$\,K over a grid of radiative transfer models of
embedded cores, two independent determinations of the same quantity that agree
to within the width of the distributions.

\subsection{Radiative transfer models}
\label{sec:models}

The tests above establish that the correction acts where it should and not
elsewhere, but they cannot show that the corrected values are closer to the
truth, because the truth is unknown in an observed image. For that we use
models of embedded dense cores whose surface density is known exactly and
whose emission is not isothermal along the line of sight.

Each model is a truncated isothermal Bonnor--Ebert sphere
\citep{Ebert1955, Bonnor1956} of gas mass $M$ and projected half-maximum
width $H$, embedded at the centre of a uniform spherical cloud of surface
density $\mathcal{D}_{\rm cl}$, whose radius is thirty times that of the sphere.
The sphere is specified by the three quantities $(\mathcal{D}_{\rm cl}, M, H)$,
which are directly comparable with what an observation measures, and its
dimensionless truncation radius, central density, and gas temperature follow
from them by inverting the Lane--Emden equation. The gas temperature of the
sphere is solved self-consistently against the total surface density along its
central line of sight, so that a more deeply embedded core is colder, as
radiative transfer requires. The models are described in detail in
Appendix~\ref{app:models}.

The dust temperature structure of each model is computed with the
Monte Carlo radiative transfer code \emph{radmc-3d} \citep{Dullemond+2012},
with the sphere and its cloud heated by the interstellar radiation field, on a
spherical grid of \draftnote{$N$} logarithmically spaced shells and with
\draftnote{$N_{\rm ph}$} photon packages. The dust opacity is a power law,
$\kappa_\nu = \kappa_0\,(\nu/\nu_0)^\beta$ with
$\kappa_0 = 10$\,cm$^2$\,g$^{-1}$ of dust at $\nu_0 = 1000$\,GHz and
$\beta = 2$, and a dust-to-gas mass ratio $\eta = 0.01$
\citep{Hildebrand1983, Ossenkopf+Henning1994}, the same values that are
assumed when the surface density image is derived, so that the comparison
tests the method and not the opacity.

For each model the synthetic \emph{Herschel} images of all wavebands were
computed, the surface density image at $13\farcs5$ was produced from them
exactly as for an observed field, and the correction of
Eq.~(\ref{eq:correction}) was applied to it. Figure~\ref{fig:validation}
compares the peak surface density recovered before and after the correction
with the true value, for 30 models spanning a factor of 40 in
peak-to-background contrast. The median recovered fraction rises from 0.63 to
0.84; over the models of contrast below 5, which is the range in which most
observed cores lie, it rises from 0.69 to 0.85.

The models of the highest contrast remain under-corrected, at 0.66 against
0.53 before the correction. They are the most compact objects of the sample,
and they illustrate the limitation of Sect.~\ref{sec:limits}: a structure
whose width approaches the angular resolution is mixed by the beam with its
warmer surroundings, so the temperature measured for its level is too high.

\begin{figure}
  \centering
  \includegraphics[width=\linewidth]{fig_model_validation.pdf}
  \caption{Peak surface density recovered from synthetic \emph{Herschel}
  images of embedded Bonnor--Ebert spheres, divided by the true value
  convolved to the same resolution, before (open circles) and after (filled
  squares) the correction of Eq.~(\ref{eq:correction}), as a function of the
  peak-to-background contrast of the model. The grey line marks exact
  recovery.}
  \label{fig:validation}
\end{figure}

% =====================================================================
\section{Limitations}
\label{sec:limits}

Three limitations should be stated plainly.

Structures narrower than about twice the angular resolution are
under-corrected. The beam mixes such a structure with its warmer surroundings
across the sky, so the temperature measured for its level is too high and the
factor of Eq.~(\ref{eq:correction}) too small. The effect is visible in
Fig.~\ref{fig:leveltemp} as the flattening of the temperature towards the
narrowest levels. It is intrinsic to a correction driven by a measured
temperature, and a correction of the photometry of individual sources, which
can use their measured sizes explicitly, remains the more accurate route for
that purpose.

Internally heated sources are excluded rather than corrected
(Sect.~\ref{sec:protostars}). Representing them would require more than one
temperature per level, which is a substantially more elaborate method for a
small minority of the sources.

The correction is not a deconvolution. It removes the bias of the
single-temperature assumption and leaves the beam dilution of a compact
structure untouched.

% =====================================================================
\section{Conclusions}
\label{sec:conclusions}

We have described a method that removes the dust temperature bias from derived
surface density images of molecular clouds.

\begin{enumerate}
\item An image is decomposed into levels by a nested set of morphological
      openings. The levels are non-negative, they sum exactly to the image, and
      each contains the material within a range of widths
      (Eqs.~\ref{eq:levels} and \ref{eq:sum}).
\item Each level is converted into surface density with a dust temperature
      measured for that level alone, instead of the single temperature of the
      whole line of sight. The correction is the ratio of the Planck functions
      at the two temperatures (Eq.~\ref{eq:correction}).
\item The correction is exactly unity where the line of sight is isothermal. It
      requires no extraction of sources, no assumption about their shapes, and
      no calibration on models, and it can neither create a depression nor
      deform a structure.
\item Applied to a sub-field of the Aquila cloud, it leaves the diffuse
      emission unchanged and raises the peak surface densities of the dense
      structures by a median factor of 1.7, with dust temperatures of the levels
      between 8.7 and 10.4\,K against 15\,K from the single-temperature fit.
\end{enumerate}

Masses of dense cores and linear densities of filaments are derived from
surface density images and inherit their bias. The corrected images are the
natural starting point for those measurements, which are the subject of
separate work.

% =====================================================================
\appendix

\section{The grid of embedded Bonnor--Ebert spheres}
\label{app:models}

\draftnote{To be completed once the grid computed for this paper is
available. It should give: the axes of the grid and their ranges; the
inversion of the Lane--Emden equation that turns
$(\mathcal{D}_{\rm cl}, M, H)$ into $(\xi_{\rm max}, \rho_{\rm c}, T)$; the
self-consistent solution for the gas temperature; the constraint that the
density at the outer boundary of the sphere must not fall below the density of
the cloud, which sets the lowest core mass at each cloud surface density; and
the \emph{radmc-3d} setup, with the radiation field, the grid, and the number
of photon packages.}

\begin{acknowledgements}
\draftnote{To be written.}
\end{acknowledgements}

\bibliographystyle{aa}
\bibliography{hires_correction}

\end{document}
